\subsection{从错误输出到完整攻击链：模型等级、误用与系统后果}

slides 20--30 用逐步 walkthrough 说明哪里会出错。攻击者可控制用户输入、外部内容、工具返回或 memory；模型可能产生不准确、不安全或被操纵的输出；host 又可能把输出当作文本、代码、数据库参数或下游 prompt。模型安全等级从理想 L0 到易受 prompt engineering、jailbreak 与 adversarial input 影响的现实等级。Misuse 既可能直接滥用模型，也可能借 agentic system 攻击外部系统。

\lecturefigure{slides-images/slide-020.png}{官方 slide 20：Agentic Hybrid System 的失败 walkthrough 起点}
\lecturefigure{slides-images/slide-022.png}{官方 slide 22：外部输入与模型决策进入攻击链}
\lecturefigure{slides-images/slide-024.png}{官方 slide 24：工具与环境扩大后果}
\lecturefigure{slides-images/slide-026.png}{官方 slide 26：完整的 what-could-go-wrong 链路}
\lecturefigure{slides-images/slide-027.png}{官方 slide 27：LLM output 可成为攻击链的一部分}
\lecturefigure{slides-images/slide-028.png}{官方 slide 28：Model security levels}
\lecturefigure{slides-images/slide-029.png}{官方 slide 29：Model misuse 与 system misuse}
\lecturefigure{slides-images/slide-030.png}{官方 slide 30：Agentic Systems 中的代表性攻击类型}

\begin{knowledgebox}{读图：危险来自类型转换}
同一串模型输出若只显示给用户，最多是内容风险；若拼入 SQL、shell、HTML、tool argument 或另一个高权限 Agent 的 prompt，语义就变成可执行控制。安全设计应在每次类型转换处重新验证，而不是认为“前面已经做过一次内容过滤”。
\end{knowledgebox}

\begin{warningbox}{输出可信度不能由语言流畅度推断}
模型可以以高置信语气生成语法正确但危险的命令。Host 必须把模型视为 untrusted planner：参数化查询、typed schema、allowlist、dry-run、transaction 和审批都应由确定性层实施，不能由模型自我声明“这个动作安全”。
\end{warningbox}

